\subsection{Introducing the Logarithm} 
\begin{definition}[Logarithm]
We define the \term{logarithm base $b$ of $x$} (written as $\log_b(x)$) as the function such that if
\[
f(x)=b^x
\]
then
\[
f\inv(x)=\log_b(x)
\]
The number $b$ is called the \term{base}, and the number $x$ to which the log is applied is called the \term{argument}.

Logarithm functions are defined for $b\makeblank[1in],b\makeblank[1in]$, just like \makeblank[1in] functions. 
\end{definition}
This means that statement involving an exponent can be rewritten in in the form of a statement about a logarithm.
\begin{Example}
Rewrite the statement $5^3=125$ in log form.
\begin{lpic}{}{1}
\end{lpic}
Since the original exponential statement is \makeblank[1in], the log form statement is also \makeblank[1in].
\end{Example}
We can also rewrite a statement about logarithms in exponential form.
\begin{Example}
Rewrite the statement $\log_3(81)=4$ in exponential form, and determine whether it's true.
\begin{lpic}{}{1}
\end{lpic}
Since the exponential statement is \makeblank[1in], the log form statement is also \makeblank[1in].
\end{Example}
Like other functions, we can compose a logarithm with something \makeblank[1in] and still have a logarithm function. For example, the following are logarithm functions:
\[
f(x)=\log_4(x-2)
\qquad
g(x)=-2\log_{1/2}(x)+4
\qquad
h(x)=3\log_2\left(\frac{1}{3}x-5\right)-2
\]
\begin{Note}
If the argument is a complicated expression, we write it in parentheses. If the argument contains a single term, the parentheses are optional.
\end{Note}